\section{Subgroup descent in the intermediate range}
\label{sec:subgroup-descent}

Fix, for this section, a sequence \(S\) satisfying
\eqref{eq:ambient-length} and \eqref{eq:middle-range}; retain its
reflection and rotation counts \(a,b\). For a sequence \(Y\) over
\(G\) and a subgroup \(K\le A\), let \(T_K(Y)\) be the sequence over
\(\Dih(A/K)\) obtained by taking all reflections of \(Y\), together
with all rotations whose vectors lie outside \(K\), and projecting their
vectors modulo \(K\). Rotations with vectors in \(K\) are omitted.
Selections in \(T_K(Y)\) always refer to the corresponding positions of
\(Y\).

We prove two statements simultaneously. For every proper subgroup
\(K<A\), let \(\mathcal P_S(K)\) denote the assertion
\[
  m_K(S)\ge b-[A:K]+2
  \quad\Longrightarrow\quad
  S\text{ has a product-one subsequence of length }2Q.
\]
Let \(\mathcal Q(K)\) denote the following universal assertion: every
sequence \(Y\) over \(G\) of length \(2Q+D_A\), with the same counts
\(a,b\), has an all-rotation product-one subsequence of length \(2Q\)
provided that
\begin{enumerate}
\item \(m_K(Y)\ge b-[A:K]+2\), and
\item \(T_K(Y)\) has no product-one subsequence containing a reflection.
\end{enumerate}
The all-rotation conclusion is essential because the proof of
\(\mathcal Q(K)\) translates rotation vectors before descending to a
smaller subgroup.

\begin{theorem}[Subgroup descent]\label{thm:subgroup-descent}
For every proper subgroup \(K<A\), both \(\mathcal P_S(K)\) and
\(\mathcal Q(K)\) hold.
\end{theorem}

\subsection{Three subgroup estimates}

\begin{lemma}[Davenport inequality]\label{lem:Davenport-inequality}
For every subgroup \(K\le A\),
\begin{equation}\label{eq:Davenport-inequality}
  D_A\ge \D(K)+\D(A/K)-1.
\end{equation}
\end{lemma}

\begin{proof}
Take a zero-sum-free sequence in \(K\) of length \(\D(K)-1\) and a
zero-sum-free sequence in \(A/K\) of length \(\D(A/K)-1\), and lift
the latter arbitrarily to \(A\). A zero-sum subsequence of the
concatenation would project to a zero-sum subsequence of the quotient
sequence. It therefore uses no lifted quotient term, after which it
would be a zero-sum subsequence of the sequence in \(K\). Thus the
concatenation is zero-sum-free, giving \eqref{eq:Davenport-inequality}.
\end{proof}

\begin{lemma}[Target-length estimate]\label{lem:target-length}
Let \(0<K<A\), put \(I=[A:K]\) and \(h=|K|\), so that \(Q=Ih\).
Then
\begin{equation}\label{eq:target-length}
  2Q-D_A-I+2\ge Q-I+2\ge h.
\end{equation}
\end{lemma}

\begin{proof}
The first inequality is \(D_A\le Q\). For the second,
\[
  Q-I+2-h=Ih-I-h+2=(I-1)(h-1)+1>0.
\]
\end{proof}

\begin{lemma}[Transitivity of concentration]\label{lem:concentration-transitivity}
Let \(H<K<A\), put \(I=[A:K]\) and \(J=[K:H]\), and let \(Y\) be a
sequence with
\[
  m_K(Y)\ge b-I+2.
\]
If at least \(m_K(Y)-J+2\) of the rotations counted by \(m_K(Y)\)
actually lie in \(H\), then
\[
  m_H(Y)\ge b-[A:H]+2.
\]
\end{lemma}

\begin{proof}
Since \([A:H]=IJ\),
\[
  m_H(Y)\ge b-I-J+4\ge b-IJ+2,
\]
because \(IJ-I-J+2=(I-1)(J-1)+1>0\).
\end{proof}

\begin{lemma}[Correction of a quotient defect]\label{lem:defect-correction}
Let \(K\le A\). Suppose a selected sequence \(U\) has a positive even
number of reflections and admits an ordering whose product has
\(A\)-coordinate \(z\in K\). Let \(C_0\) be a disjoint rotation
sequence, with prescribed signs, whose signed vector sum is \(-z\).
Then \(UC_0\) is product one.
\end{lemma}

\begin{proof}
Read from the fixed ordering of \(U\) the sign of each vector. The
reflection signs are balanced, because they alternate and their number
is positive and even. Together with the prescribed signs on \(C_0\),
the total signed vector sum is zero. Apply
\Cref{lem:sign-realization} to the entire selected sequence.
\end{proof}

\subsection{A reflection-containing quotient block}
\label{subsec:reflection-quotient}

Let \(0<K<A\), suppose the hypothesis of \(\mathcal P_S(K)\) holds,
and let \(C\) be the sequence of the \(M=m_K(S)\) rotations of \(S\)
whose vectors lie in \(K\). Let \(B\) be the sequence of the remaining
\(c=b-M\) rotations, and put \(T=T_K(S)\).

\begin{proposition}\label{prop:reflection-quotient}
If \(T\) has a product-one subsequence containing a reflection, then
either \(S\) has a product-one subsequence of length \(2Q\), or there
is a strict subgroup \(H<K\) such that
\[
  m_H(S)\ge b-[A:H]+2.
\]
\end{proposition}

\begin{proof}
Start with a reflection-containing product-one subsequence of \(T\).
Continue removing nonempty product-one subsequences from the current
remainder until the final remainder \(R\) is product-one-free. Let
\(U\) be the union of all removed blocks. Concatenating one product-one
ordering for each block makes \(U\) product one in \(\Dih(A/K)\), and
\[
  |U|=a+c-|R|.
\]
The sequence \(U\) contains a positive even number of reflections. By
\eqref{eq:small-Davenport-dihedral}, applied to \(A/K\),
\[
  |R|\le \D(A/K).
\]
Set
\[
  \tau=D_A-|R|.
\]
By \Cref{lem:Davenport-inequality} and
\Cref{prop:plus-minus-bound},
\begin{equation}\label{eq:reflection-surplus}
  \tau\ge \D(K)-1\ge \Dpm(K)-1.
\end{equation}

Lift the fixed quotient ordering of \(U\) to the corresponding terms of
\(S\). Its product has \(C_2\)-coordinate zero and is trivial modulo
\(K\), so its \(A\)-coordinate is a defect \(z\in K\). Put
\[
  m=M-\tau=M-D_A+|R|.
\]
The concentration hypothesis and \Cref{lem:target-length} give
\begin{align*}
  m&\ge M-D_A\\
   &\ge b-[A:K]+2-D_A\\
   &=2Q-a-[A:K]+2\\
   &\ge 2Q-D_A-[A:K]+2\\
   &\ge |K|.
\end{align*}
Equation \eqref{eq:reflection-surplus} gives
\(M\ge m+\Dpm(K)-1\). Thus the structural part of
\Cref{thm:GMO} applies to \(C\) in the ambient group \(K\), with
weight set \(\{-1,1\}\) and target \(m\).

If \(\Sigma_m^{\{-1,1\}}(C)=K\), choose \(m\) rotations of \(C\),
with signs, whose signed sum is \(-z\). Their positions are disjoint
from those of \(U\), because \(T\) contains no rotation from \(K\).
By \Cref{lem:defect-correction}, their union with the lift of \(U\) is
product one. Its length is
\begin{align*}
  |U|+m
  &=a+c-|R|+M-D_A+|R|\\
  &=a+b-D_A=2Q.
\end{align*}

If the spectrum is not all of \(K\), the structural alternative gives
a strict subgroup \(H<K\) and at least \(M-[K:H]+2\) terms of \(C\)
satisfying \eqref{eq:GMO-weight-coset}. For each such vector \(x\), the
elements \(x\) and \(-x\) lie in a common \(H\)-coset, so \(2x\in H\).
The quotient \(K/H\) has odd order, and hence \(x\in H\). Now
\Cref{lem:concentration-transitivity} gives
\(m_H(S)\ge b-[A:H]+2\).
\end{proof}

\subsection{The rotation-only quotient case}
\label{subsec:rotation-quotient}

For \(\alpha\in A\) and a sequence \(Y\) over \(G\), define
\(\tau_\alpha(Y)\) by replacing every rotation \((v,0)\) with
\((v-\alpha,0)\), leaving every reflection unchanged. This is a
transformation of the values attached to the positions; it is not a
group homomorphism.

\begin{lemma}[Persistence of the quotient obstruction]
\label{lem:quotient-obstruction}
Let \(H<K\le A\) and \(\alpha\in K\). If \(T_K(Y)\) has no
product-one subsequence containing a reflection, then
\(T_H(\tau_\alpha(Y))\) has no such subsequence.
\end{lemma}

\begin{proof}
Suppose that a reflection-containing product-one subsequence existed in
\(T_H(\tau_\alpha(Y))\), and fix a product-one ordering. Project this
ordered block further modulo \(K\). A selected rotation whose original
vector \(v\) lies in \(K\) has translated vector \(v-\alpha\in K\), so
it becomes the identity modulo \(K\); delete all such identity rotations
from the ordering. A selected rotation whose original vector lies
outside \(K\) has the same image modulo \(K\) before and after the
translation. Reflections are unchanged. The remaining ordered block is
a reflection-containing product-one subsequence of \(T_K(Y)\), a
contradiction.
\end{proof}

\begin{proposition}\label{prop:rotation-quotient}
Let \(0<K<A\). Suppose that \(Y\) has length \(2Q+D_A\), has the same
counts \(a,b\) as \(S\), satisfies
\[
  m_K(Y)\ge b-[A:K]+2,
\]
and \(T_K(Y)\) has no product-one subsequence containing a reflection.
Then either
\begin{enumerate}
\item \(Y\) has an all-rotation product-one subsequence of length
      \(2Q\), or
\item there are \(H<K\) and \(\alpha\in K\) such that
      \(\tau_\alpha(Y)\) satisfies the hypotheses of \(\mathcal Q(H)\).
\end{enumerate}
\end{proposition}

\begin{proof}
Let \(C\) be the sequence of the \(M=m_K(Y)\) rotations in \(K\), and
let \(B\) be the sequence of the remaining \(c=b-M\) rotations. In the
quotient \(A/K\), successively remove nonempty zero-sum subsequences
from the image of \(B\) until a zero-sum-free remainder \(B_0\)
remains. Let \(B'=BB_0^{-1}\), using the corresponding original
positions, and put \(r_0=|B_0|\). The vector sum of \(B'\) is an
element
\begin{equation}\label{eq:rotation-defect}
  z=\sigma(B')\in K.
\end{equation}

The quotient sequence formed by all \(a\) reflections and the terms of
\(B_0\) is product-one-free. Indeed, a rotation-only product-one
subsequence would give a nonempty zero-sum subsequence of \(B_0\), and
a product-one subsequence containing a reflection is excluded by the
hypothesis on \(T_K(Y)\). Hence, by
\eqref{eq:small-Davenport-dihedral},
\begin{equation}\label{eq:rotation-remainder}
  a+r_0\le \D(A/K).
\end{equation}
Set
\[
  d_0=D_A-a-r_0.
\]
Using \Cref{lem:Davenport-inequality} and
\eqref{eq:rotation-remainder},
\begin{equation}\label{eq:rotation-surplus}
  d_0\ge D_A-\D(A/K)\ge \D(K)-1.
\end{equation}

Put \(m=M-d_0\). The concentration hypothesis gives
\begin{align*}
  m&\ge b-[A:K]+2-(D_A-a-r_0)\\
   &=2Q-[A:K]+r_0+2\\
   &\ge 2Q-[A:K]+2\ge |K|,
\end{align*}
where the last inequality follows from
\Cref{lem:target-length}. Equation \eqref{eq:rotation-surplus} gives
\(M\ge m+\D(K)-1\). Apply the structural part of
\Cref{thm:GMO} to \(C\) in \(K\), with weight set \(\{1\}\) and
target \(m\).

Suppose first that \(\Sigma_m^{\{1\}}(C)=K\). Positional
complementation inside \(C\) gives
\[
  \Sigma_{d_0}^{\{1\}}(C)
  =\sigma(C)-\Sigma_m^{\{1\}}(C)=K.
\]
Choose a \(d_0\)-term subsequence \(D_0\mid C\) with
\(\sigma(D_0)=\sigma(C)+z\). The rotation sequences
\(CD_0^{-1}\) and \(B'\) have disjoint positions and total vector sum
\[
  \sigma(C)-\sigma(D_0)+\sigma(B')=-z+z=0.
\]
Their total length is
\begin{align*}
  |CD_0^{-1}|+|B'|
  &=M-d_0+c-r_0\\
  &=b-(D_A-a-r_0)-r_0=2Q.
\end{align*}
This proves the first alternative.

If the spectrum is not all of \(K\), the structural alternative gives
a strict subgroup \(H<K\), an element \(\alpha\in K\), and at least
\(M-[K:H]+2\) terms of \(C\) whose vectors lie in \(\alpha+H\).
After applying \(\tau_\alpha\), these rotations lie in \(H\), and
\Cref{lem:concentration-transitivity} gives
\[
  m_H(\tau_\alpha(Y))\ge b-[A:H]+2.
\]
The translation preserves the length and the counts \(a,b\), while
\Cref{lem:quotient-obstruction} preserves the absence of a
reflection-containing quotient block. Thus \(\tau_\alpha(Y)\) satisfies
the hypotheses of \(\mathcal Q(H)\).
\end{proof}

\begin{lemma}[Pullback through a rotation translation]
\label{lem:translation-pullback}
If the positions of an all-rotation subsequence of length \(2Q\) form a
product-one sequence in \(\tau_\alpha(Y)\), then the same positions form
a product-one sequence in \(Y\).
\end{lemma}

\begin{proof}
Pulling the selected rotations back adds \(\alpha\) to each vector, so
the total vector sum changes by \(2Q\alpha=0\), using
\(\exp(A)\mid Q\). There are no reflections, so no alternating signs
occur.
\end{proof}

\subsection{Simultaneous induction}

\begin{proof}[Proof of \Cref{thm:subgroup-descent}]
We use strong induction on the order of the subgroup.

For \(K=0\), the hypothesis of \(\mathcal P_S(0)\) says that the
number \(m_0\) of identity rotations satisfies
\[
  m_0\ge b-Q+2=Q+D_A-a+2\ge D_A.
\]
Thus \(\mathcal P_S(0)\) follows from
\Cref{lem:identity-padding}.

To prove \(\mathcal Q(0)\), let \(Y\) satisfy its hypotheses, and let
\(m_0\) again be the number of identity rotations. The sequence
\(T_0(Y)\) consists of all reflections and all nonidentity rotations.
Successively remove nonempty product-one subsequences from \(T_0(Y)\)
until a product-one-free remainder remains, and let \(P\) be the union
of the removed blocks. The quotient obstruction in the definition of
\(\mathcal Q(0)\) forces every removed block to consist only of
rotations. Hence \(P\) is an all-rotation product-one sequence. The
terminal remainder has length at most \(D_A\), so
\[
  |P|\ge (2Q+D_A-m_0)-D_A=2Q-m_0.
\]
Moreover, the concentration hypothesis gives
\(m_0\ge b-Q+2\), and therefore
\[
  |P|\le b-m_0\le Q-2<2Q.
\]
There are enough unused identity rotations to extend \(P\) to an
all-rotation product-one subsequence of length \(2Q\). This proves
\(\mathcal Q(0)\).

Now let \(0<K<A\), and assume both statements are known for every
strict subgroup of \(K\). To prove \(\mathcal P_S(K)\), consider
\(T_K(S)\). If it has a product-one subsequence containing a reflection,
\Cref{prop:reflection-quotient} either gives the required subsequence
immediately or gives a strict subgroup \(H<K\) satisfying the hypothesis
of \(\mathcal P_S(H)\). The induction hypothesis applies in the latter
case. If \(T_K(S)\) has no such subsequence,
\Cref{prop:rotation-quotient} either gives an all-rotation subsequence of
length \(2Q\) directly or gives \(H<K\) and \(\alpha\in K\) such that
\(\tau_\alpha(S)\) satisfies the hypotheses of \(\mathcal Q(H)\). The
induction hypothesis gives the required all-rotation subsequence after
translation, and \Cref{lem:translation-pullback} transfers it back to
\(S\).

Finally, let \(Y\) satisfy the hypotheses of \(\mathcal Q(K)\). Its
quotient obstruction already rules out the reflection-containing case.
Apply \Cref{prop:rotation-quotient}. The first alternative is the desired
conclusion. In the second, the translated sequence satisfies
\(\mathcal Q(H)\) for a strict subgroup \(H<K\); apply the induction
hypothesis and then \Cref{lem:translation-pullback}.
\end{proof}
